% =====================================================================
% CHƯƠNG 6: KẾT LUẬN VÀ HƯỚNG PHÁT TRIỂN
% ĐỀ TÀI: XÂY DỰNG HỆ THỐNG ĐỌC NHÃN MÁC VÀ CẢNH BÁO DỊ ỨNG THỰC PHẨM
% =====================================================================

\chapter{Kết luận và hướng phát triển}
\label{ch:ketluan}

\section{Kết luận}

Đồ án ``Xây dựng hệ thống đọc nhãn mác và cảnh báo dị ứng thực
phẩm'' đã xây dựng được hệ thống Allergy Warning nhằm hỗ trợ
người dùng nhận diện thành phần trên nhãn thực phẩm và phát hiện
các dị nguyên liên quan đến hồ sơ dị ứng cá nhân.

Hệ thống được phát triển bằng Flask, sử dụng Gemini API để phân
tích hình ảnh, Cloudinary để lưu trữ ảnh, TiDB Cloud để quản lý
dữ liệu và Open Food Facts làm nguồn tham khảo thông tin sản phẩm.
Ứng dụng được triển khai trên Render với các chức năng chính:
quản lý tài khoản, hồ sơ dị ứng, quét nhãn thực phẩm, hiển thị
cảnh báo và lưu lịch sử quét.

Qua quá trình thực hiện, đồ án đã vận dụng kiến thức về phát triển
ứng dụng web, quản lý cơ sở dữ liệu và tích hợp trí tuệ nhân tạo.
Đồng thời, quá trình xây dựng cho thấy chất lượng kết quả phụ thuộc
vào hình ảnh đầu vào, dữ liệu dị nguyên và khả năng phân tích của
dịch vụ AI.

Mặc dù đã hoàn thành các chức năng chính, hệ thống vẫn có thể
nhận diện thiếu hoặc sai thông tin trong một số trường hợp.
Do đó, Allergy Warning hiện đóng vai trò hỗ trợ tham khảo, không
thay thế việc kiểm tra nhãn gốc hoặc tư vấn chuyên môn.

\section{Bài học rút ra}

Trong quá trình thực hiện đồ án, có thể rút ra một số bài học:

\begin{itemize}
	\item Cần kiểm tra chất lượng hình ảnh đầu vào để hạn chế
	sai lệch trong quá trình nhận diện.
	
	\item Kết quả từ AI cần được kiểm tra và xử lý trước khi
	sử dụng để đưa ra cảnh báo.
	
	\item Dữ liệu dị nguyên cần được chuẩn hóa, đặc biệt là
	các tên gọi tương đương và những thành phần dễ nhầm lẫn.
	
	\item Cần kiểm thử toàn bộ quy trình, từ tải ảnh, phân tích,
	đối chiếu đến hiển thị kết quả và lưu lịch sử.
\end{itemize}

\section{Hướng phát triển}

Trong thời gian tới, hệ thống có thể được hoàn thiện theo các
hướng sau:

\begin{itemize}
	\item Mở rộng dữ liệu:
	Bổ sung dị nguyên và các tên gọi tương đương để nâng cao
	khả năng nhận diện, đối chiếu thành phần.
	
	\item Cải thiện xử lý hình ảnh:
	Bổ sung kiểm tra chất lượng ảnh và các bước tiền xử lý
	nhằm hỗ trợ nhận diện nội dung nhãn chính xác hơn.
	
	\item Hoàn thiện phân tích và cảnh báo:
	Cải thiện quy tắc chuẩn hóa, phân biệt thành phần nguyên liệu
	với cảnh báo có thể chứa vết dị nguyên.
	
	\item Nâng cao khả năng kiểm thử:
	Xây dựng tập ảnh nhãn thực phẩm đa dạng và sử dụng
	Precision, Recall, F1-score để đánh giá khi có đủ dữ liệu.
	
	\item Tăng cường độ ổn định và bảo mật:
	Hoàn thiện xử lý lỗi khi các dịch vụ bên ngoài hoặc cơ sở
	dữ liệu gặp sự cố, đồng thời bảo vệ thông tin người dùng.
	
	\item Hoàn thiện giao diện:
	Cải thiện cách trình bày kết quả, cảnh báo và hướng dẫn
	để người dùng dễ theo dõi, kiểm tra thông tin.
\end{itemize}

Các hướng phát triển trên nhằm nâng cao độ tin cậy, khả năng
sử dụng và tính ổn định của hệ thống trong thực tế.